\section{Inference cost}
\label{sec:cost}

% receipt: evidence/tools/reports/paper/timing_2026-09-10/timing.runs/cq_release.complete.json
We measure synchronized native-inference latency on an NVIDIA RTX 3090,
with no competing GPU compute process detected, at a fixed 270~W power limit. A fixed development subset supplies 500 images
from 250 factor-cell scenes. Each system runs in
a fresh process with batch size 1 and native FP16 autocast. The harness excludes
file decode and 30 warm-up images, while retaining resize, normalization,
device transfer, model execution, routing, warping, test-time augmentation, host
transfer, and notation construction whenever the native pipeline uses them.

% receipt: evidence/tools/reports/paper/timing_2026-09-10/timing.runs.json
The release checkpoint has median latency 64.41~ms per image. The raw member
requires 79.69~ms, while the shipped ensemble requires 195.52~ms. The ten
adapted checkpoints lie between 71.50 and 73.12~ms, so their measured median
overhead relative to the release checkpoint is 7.09--8.71~ms in this harness.
Table~\ref{tab:timing} reports medians and ninetieth percentiles for every
system, and Figure~\ref{fig:timing} visualizes the same receipt-bound values.

\begingroup\small
\setlength{\tabcolsep}{4pt}
\begin{table}[htbp]\centering
\caption{Synchronized single-image latency on the same isolated device.}\label{tab:timing}
\begin{tabular}{lrrr}
\toprule
System & % receipt: evidence/tools/reports/paper/timing_2026-09-10/timing.complete.json; source: evidence/tools/reports/paper/timing_2026-09-10/timing.runs.json; column definitions
Images & % receipt: evidence/tools/reports/paper/timing_2026-09-10/timing.complete.json; source: evidence/tools/reports/paper/timing_2026-09-10/timing.runs.json; column definitions
Median (ms) & % receipt: evidence/tools/reports/paper/timing_2026-09-10/timing.complete.json; source: evidence/tools/reports/paper/timing_2026-09-10/timing.runs.json; column definitions
P90 (ms) \\ % receipt: evidence/tools/reports/paper/timing_2026-09-10/timing.complete.json; source: evidence/tools/reports/paper/timing_2026-09-10/timing.runs.json; column definitions
\midrule

ChessQueries release & % receipt: evidence/tools/reports/paper/timing_2026-09-10/timing.complete.json; source: evidence/tools/reports/paper/timing_2026-09-10/timing.runs.json; cq_release/summary
500 & % receipt: evidence/tools/reports/paper/timing_2026-09-10/timing.complete.json; source: evidence/tools/reports/paper/timing_2026-09-10/timing.runs.json; cq_release/summary
64.41 & % receipt: evidence/tools/reports/paper/timing_2026-09-10/timing.complete.json; source: evidence/tools/reports/paper/timing_2026-09-10/timing.runs.json; cq_release/summary
65.72 \\ % receipt: evidence/tools/reports/paper/timing_2026-09-10/timing.complete.json; source: evidence/tools/reports/paper/timing_2026-09-10/timing.runs.json; cq_release/summary

shipped ensemble & % receipt: evidence/tools/reports/paper/timing_2026-09-10/timing.complete.json; source: evidence/tools/reports/paper/timing_2026-09-10/timing.runs.json; private_full/summary
500 & % receipt: evidence/tools/reports/paper/timing_2026-09-10/timing.complete.json; source: evidence/tools/reports/paper/timing_2026-09-10/timing.runs.json; private_full/summary
195.52 & % receipt: evidence/tools/reports/paper/timing_2026-09-10/timing.complete.json; source: evidence/tools/reports/paper/timing_2026-09-10/timing.runs.json; private_full/summary
244.22 \\ % receipt: evidence/tools/reports/paper/timing_2026-09-10/timing.complete.json; source: evidence/tools/reports/paper/timing_2026-09-10/timing.runs.json; private_full/summary

raw member & % receipt: evidence/tools/reports/paper/timing_2026-09-10/timing.complete.json; source: evidence/tools/reports/paper/timing_2026-09-10/timing.runs.json; private_raw/summary
500 & % receipt: evidence/tools/reports/paper/timing_2026-09-10/timing.complete.json; source: evidence/tools/reports/paper/timing_2026-09-10/timing.runs.json; private_raw/summary
79.69 & % receipt: evidence/tools/reports/paper/timing_2026-09-10/timing.complete.json; source: evidence/tools/reports/paper/timing_2026-09-10/timing.runs.json; private_raw/summary
81.07 \\ % receipt: evidence/tools/reports/paper/timing_2026-09-10/timing.complete.json; source: evidence/tools/reports/paper/timing_2026-09-10/timing.runs.json; private_raw/summary

warp-high member & % receipt: evidence/tools/reports/paper/timing_2026-09-10/timing.complete.json; source: evidence/tools/reports/paper/timing_2026-09-10/timing.runs.json; private_warp_high/summary
500 & % receipt: evidence/tools/reports/paper/timing_2026-09-10/timing.complete.json; source: evidence/tools/reports/paper/timing_2026-09-10/timing.runs.json; private_warp_high/summary
195.23 & % receipt: evidence/tools/reports/paper/timing_2026-09-10/timing.complete.json; source: evidence/tools/reports/paper/timing_2026-09-10/timing.runs.json; private_warp_high/summary
197.55 \\ % receipt: evidence/tools/reports/paper/timing_2026-09-10/timing.complete.json; source: evidence/tools/reports/paper/timing_2026-09-10/timing.runs.json; private_warp_high/summary

warp-low member & % receipt: evidence/tools/reports/paper/timing_2026-09-10/timing.complete.json; source: evidence/tools/reports/paper/timing_2026-09-10/timing.runs.json; private_warp_low/summary
500 & % receipt: evidence/tools/reports/paper/timing_2026-09-10/timing.complete.json; source: evidence/tools/reports/paper/timing_2026-09-10/timing.runs.json; private_warp_low/summary
168.98 & % receipt: evidence/tools/reports/paper/timing_2026-09-10/timing.complete.json; source: evidence/tools/reports/paper/timing_2026-09-10/timing.runs.json; private_warp_low/summary
170.85 \\ % receipt: evidence/tools/reports/paper/timing_2026-09-10/timing.complete.json; source: evidence/tools/reports/paper/timing_2026-09-10/timing.runs.json; private_warp_low/summary

LoRA dark n500 s303 & % receipt: evidence/tools/reports/paper/timing_2026-09-10/timing.complete.json; source: evidence/tools/reports/paper/timing_2026-09-10/timing.runs.json; lora_dark_n500_s303/summary
500 & % receipt: evidence/tools/reports/paper/timing_2026-09-10/timing.complete.json; source: evidence/tools/reports/paper/timing_2026-09-10/timing.runs.json; lora_dark_n500_s303/summary
71.88 & % receipt: evidence/tools/reports/paper/timing_2026-09-10/timing.complete.json; source: evidence/tools/reports/paper/timing_2026-09-10/timing.runs.json; lora_dark_n500_s303/summary
72.64 \\ % receipt: evidence/tools/reports/paper/timing_2026-09-10/timing.complete.json; source: evidence/tools/reports/paper/timing_2026-09-10/timing.runs.json; lora_dark_n500_s303/summary

LoRA both n500 s303 & % receipt: evidence/tools/reports/paper/timing_2026-09-10/timing.complete.json; source: evidence/tools/reports/paper/timing_2026-09-10/timing.runs.json; lora_both_n500_s303/summary
500 & % receipt: evidence/tools/reports/paper/timing_2026-09-10/timing.complete.json; source: evidence/tools/reports/paper/timing_2026-09-10/timing.runs.json; lora_both_n500_s303/summary
72.40 & % receipt: evidence/tools/reports/paper/timing_2026-09-10/timing.complete.json; source: evidence/tools/reports/paper/timing_2026-09-10/timing.runs.json; lora_both_n500_s303/summary
73.32 \\ % receipt: evidence/tools/reports/paper/timing_2026-09-10/timing.complete.json; source: evidence/tools/reports/paper/timing_2026-09-10/timing.runs.json; lora_both_n500_s303/summary

LoRA dark n2000 s303 & % receipt: evidence/tools/reports/paper/timing_2026-09-10/timing.complete.json; source: evidence/tools/reports/paper/timing_2026-09-10/timing.runs.json; lora_dark_n2000_s303/summary
500 & % receipt: evidence/tools/reports/paper/timing_2026-09-10/timing.complete.json; source: evidence/tools/reports/paper/timing_2026-09-10/timing.runs.json; lora_dark_n2000_s303/summary
72.29 & % receipt: evidence/tools/reports/paper/timing_2026-09-10/timing.complete.json; source: evidence/tools/reports/paper/timing_2026-09-10/timing.runs.json; lora_dark_n2000_s303/summary
73.08 \\ % receipt: evidence/tools/reports/paper/timing_2026-09-10/timing.complete.json; source: evidence/tools/reports/paper/timing_2026-09-10/timing.runs.json; lora_dark_n2000_s303/summary

LoRA both n2000 s303 & % receipt: evidence/tools/reports/paper/timing_2026-09-10/timing.complete.json; source: evidence/tools/reports/paper/timing_2026-09-10/timing.runs.json; lora_both_n2000_s303/summary
500 & % receipt: evidence/tools/reports/paper/timing_2026-09-10/timing.complete.json; source: evidence/tools/reports/paper/timing_2026-09-10/timing.runs.json; lora_both_n2000_s303/summary
71.50 & % receipt: evidence/tools/reports/paper/timing_2026-09-10/timing.complete.json; source: evidence/tools/reports/paper/timing_2026-09-10/timing.runs.json; lora_both_n2000_s303/summary
72.50 \\ % receipt: evidence/tools/reports/paper/timing_2026-09-10/timing.complete.json; source: evidence/tools/reports/paper/timing_2026-09-10/timing.runs.json; lora_both_n2000_s303/summary

LoRA dark n2000 s304 & % receipt: evidence/tools/reports/paper/timing_2026-09-10/timing.complete.json; source: evidence/tools/reports/paper/timing_2026-09-10/timing.runs.json; lora_dark_n2000_s304/summary
500 & % receipt: evidence/tools/reports/paper/timing_2026-09-10/timing.complete.json; source: evidence/tools/reports/paper/timing_2026-09-10/timing.runs.json; lora_dark_n2000_s304/summary
71.86 & % receipt: evidence/tools/reports/paper/timing_2026-09-10/timing.complete.json; source: evidence/tools/reports/paper/timing_2026-09-10/timing.runs.json; lora_dark_n2000_s304/summary
72.95 \\ % receipt: evidence/tools/reports/paper/timing_2026-09-10/timing.complete.json; source: evidence/tools/reports/paper/timing_2026-09-10/timing.runs.json; lora_dark_n2000_s304/summary

LoRA both n2000 s304 & % receipt: evidence/tools/reports/paper/timing_2026-09-10/timing.complete.json; source: evidence/tools/reports/paper/timing_2026-09-10/timing.runs.json; lora_both_n2000_s304/summary
500 & % receipt: evidence/tools/reports/paper/timing_2026-09-10/timing.complete.json; source: evidence/tools/reports/paper/timing_2026-09-10/timing.runs.json; lora_both_n2000_s304/summary
72.80 & % receipt: evidence/tools/reports/paper/timing_2026-09-10/timing.complete.json; source: evidence/tools/reports/paper/timing_2026-09-10/timing.runs.json; lora_both_n2000_s304/summary
73.68 \\ % receipt: evidence/tools/reports/paper/timing_2026-09-10/timing.complete.json; source: evidence/tools/reports/paper/timing_2026-09-10/timing.runs.json; lora_both_n2000_s304/summary

LoRA dark n2000 s305 & % receipt: evidence/tools/reports/paper/timing_2026-09-10/timing.complete.json; source: evidence/tools/reports/paper/timing_2026-09-10/timing.runs.json; lora_dark_n2000_s305/summary
500 & % receipt: evidence/tools/reports/paper/timing_2026-09-10/timing.complete.json; source: evidence/tools/reports/paper/timing_2026-09-10/timing.runs.json; lora_dark_n2000_s305/summary
72.50 & % receipt: evidence/tools/reports/paper/timing_2026-09-10/timing.complete.json; source: evidence/tools/reports/paper/timing_2026-09-10/timing.runs.json; lora_dark_n2000_s305/summary
73.48 \\ % receipt: evidence/tools/reports/paper/timing_2026-09-10/timing.complete.json; source: evidence/tools/reports/paper/timing_2026-09-10/timing.runs.json; lora_dark_n2000_s305/summary

LoRA both n2000 s305 & % receipt: evidence/tools/reports/paper/timing_2026-09-10/timing.complete.json; source: evidence/tools/reports/paper/timing_2026-09-10/timing.runs.json; lora_both_n2000_s305/summary
500 & % receipt: evidence/tools/reports/paper/timing_2026-09-10/timing.complete.json; source: evidence/tools/reports/paper/timing_2026-09-10/timing.runs.json; lora_both_n2000_s305/summary
73.12 & % receipt: evidence/tools/reports/paper/timing_2026-09-10/timing.complete.json; source: evidence/tools/reports/paper/timing_2026-09-10/timing.runs.json; lora_both_n2000_s305/summary
74.03 \\ % receipt: evidence/tools/reports/paper/timing_2026-09-10/timing.complete.json; source: evidence/tools/reports/paper/timing_2026-09-10/timing.runs.json; lora_both_n2000_s305/summary

LoRA dark n8000 s303 & % receipt: evidence/tools/reports/paper/timing_2026-09-10/timing.complete.json; source: evidence/tools/reports/paper/timing_2026-09-10/timing.runs.json; lora_dark_n8000_s303/summary
500 & % receipt: evidence/tools/reports/paper/timing_2026-09-10/timing.complete.json; source: evidence/tools/reports/paper/timing_2026-09-10/timing.runs.json; lora_dark_n8000_s303/summary
72.40 & % receipt: evidence/tools/reports/paper/timing_2026-09-10/timing.complete.json; source: evidence/tools/reports/paper/timing_2026-09-10/timing.runs.json; lora_dark_n8000_s303/summary
73.26 \\ % receipt: evidence/tools/reports/paper/timing_2026-09-10/timing.complete.json; source: evidence/tools/reports/paper/timing_2026-09-10/timing.runs.json; lora_dark_n8000_s303/summary

LoRA both n8000 s303 & % receipt: evidence/tools/reports/paper/timing_2026-09-10/timing.complete.json; source: evidence/tools/reports/paper/timing_2026-09-10/timing.runs.json; lora_both_n8000_s303/summary
500 & % receipt: evidence/tools/reports/paper/timing_2026-09-10/timing.complete.json; source: evidence/tools/reports/paper/timing_2026-09-10/timing.runs.json; lora_both_n8000_s303/summary
72.66 & % receipt: evidence/tools/reports/paper/timing_2026-09-10/timing.complete.json; source: evidence/tools/reports/paper/timing_2026-09-10/timing.runs.json; lora_both_n8000_s303/summary
73.57 \\ % receipt: evidence/tools/reports/paper/timing_2026-09-10/timing.complete.json; source: evidence/tools/reports/paper/timing_2026-09-10/timing.runs.json; lora_both_n8000_s303/summary

\bottomrule
\end{tabular}
\end{table}
\endgroup


These measurements compare complete native inference paths rather than equal
floating-point workloads. In particular, the ensemble and its members differ
in routing, member composition, warping, and test-time augmentation. The timing
experiment supports no accuracy claim and should not be extrapolated beyond the
recorded hardware, software, and input pipeline. Each model has one timed process;
percentiles describe variation across input images, not run-to-run repeatability.
No repeated-process confidence interval or whole-project rendering/training cost
is estimated by this table.
